\documentclass[10pt,a4paper]{article}
\usepackage[margin=2.2cm]{geometry}
\usepackage{amsmath,amssymb,amsthm,booktabs,graphicx,microtype,siunitx}
\usepackage[table]{xcolor}
\usepackage[round,authoryear]{natbib}
\usepackage[hidelinks]{hyperref}
\usepackage{caption}
\usepackage{times}
\captionsetup{font=small,labelfont=bf}
\sisetup{detect-all}
\newtheorem{proposition}{Proposition}
\newcommand{\zv}{\mathbf{z}}
\newcommand{\V}[1]{\csname v#1\endcsname}
\expandafter\def\csname vnruns\endcsname{132}
\expandafter\def\csname vgpuh\endcsname{204}
\expandafter\def\csname vspread64\endcsname{0.250}
\expandafter\def\csname vseedsd64\endcsname{0.238}
\expandafter\def\csname vanovaF\endcsname{0.954}
\expandafter\def\csname vanovaP\endcsname{0.450}
\expandafter\def\csname vanovaEta\endcsname{0.132}
\expandafter\def\csname vw64uniformd\endcsname{+0.048}
\expandafter\def\csname vw64uniformtost\endcsname{0.0142}
\expandafter\def\csname vw64gapd\endcsname{+0.080}
\expandafter\def\csname vw64gaptost\endcsname{0.0188}
\expandafter\def\csname vw64shuffledd\endcsname{-0.073}
\expandafter\def\csname vw64shuffledtost\endcsname{0.0464}
\expandafter\def\csname vw64frozenrandomd\endcsname{+0.177}
\expandafter\def\csname vw64frozenrandomtost\endcsname{0.0191}
\expandafter\def\csname vfrzVgap\endcsname{+0.097}
\expandafter\def\csname vfrzVgapP\endcsname{0.443}
\expandafter\def\csname vfrzVgapT\endcsname{0.0089}
\expandafter\def\csname vgapEntMin\endcsname{0.9867}
\expandafter\def\csname vgapEntMean\endcsname{0.9885}
\expandafter\def\csname vpowsd\endcsname{0.257}
\expandafter\def\csname vmde\endcsname{0.270}
\expandafter\def\csname vpwA\endcsname{48.5}
\expandafter\def\csname vpwB\endcsname{96.3}
\expandafter\def\csname vpwC\endcsname{100.0}
\expandafter\def\csname veffCA\endcsname{+0.532}
\expandafter\def\csname veffCAlo\endcsname{+0.292}
\expandafter\def\csname veffCAhi\endcsname{+0.771}
\expandafter\def\csname veffCAp\endcsname{=0.0023}
\expandafter\def\csname vbankCA\endcsname{-0.223}
\expandafter\def\csname vbankCAT\endcsname{0.0251}
\expandafter\def\csname vroutCA\endcsname{+0.755}
\expandafter\def\csname vlossCA\endcsname{0.1643}
\expandafter\def\csname veffCB\endcsname{+1.153}
\expandafter\def\csname veffCBlo\endcsname{+0.940}
\expandafter\def\csname veffCBhi\endcsname{+1.367}
\expandafter\def\csname veffCBp\endcsname{<10^{-4}}
\expandafter\def\csname vbankCB\endcsname{+0.175}
\expandafter\def\csname vbankCBT\endcsname{0.0242}
\expandafter\def\csname vroutCB\endcsname{+0.978}
\expandafter\def\csname vlossCB\endcsname{0.0170}
\expandafter\def\csname veffCC\endcsname{+0.080}
\expandafter\def\csname veffCClo\endcsname{-0.305}
\expandafter\def\csname veffCChi\endcsname{+0.465}
\expandafter\def\csname veffCCp\endcsname{=0.6158}
\expandafter\def\csname vbankCC\endcsname{+0.048}
\expandafter\def\csname vbankCCT\endcsname{0.0142}
\expandafter\def\csname vroutCC\endcsname{+0.032}
\expandafter\def\csname vlossCC\endcsname{0.0080}
\expandafter\def\csname veffTA\endcsname{+2.185}
\expandafter\def\csname veffTAlo\endcsname{+1.781}
\expandafter\def\csname veffTAhi\endcsname{+2.589}
\expandafter\def\csname veffTAp\endcsname{<10^{-4}}
\expandafter\def\csname vbankTA\endcsname{+0.107}
\expandafter\def\csname vbankTAT\endcsname{0.0218}
\expandafter\def\csname vroutTA\endcsname{+2.078}
\expandafter\def\csname vlossTA\endcsname{0.1194}
\expandafter\def\csname veffTB\endcsname{+2.282}
\expandafter\def\csname veffTBlo\endcsname{+1.952}
\expandafter\def\csname veffTBhi\endcsname{+2.611}
\expandafter\def\csname veffTBp\endcsname{<10^{-4}}
\expandafter\def\csname vbankTB\endcsname{+0.053}
\expandafter\def\csname vbankTBT\endcsname{0.0157}
\expandafter\def\csname vroutTB\endcsname{+2.228}
\expandafter\def\csname vlossTB\endcsname{0.0210}
\expandafter\def\csname vmMA\endcsname{+4.345}
\expandafter\def\csname vmMAlo\endcsname{+4.103}
\expandafter\def\csname vmMAhi\endcsname{+4.587}
\expandafter\def\csname vtMAg\endcsname{33.7}
\expandafter\def\csname vtMAs\endcsname{8.6}
\expandafter\def\csname vmMB\endcsname{+2.313}
\expandafter\def\csname vmMBlo\endcsname{+1.870}
\expandafter\def\csname vmMBhi\endcsname{+2.757}
\expandafter\def\csname vtMBg\endcsname{39.2}
\expandafter\def\csname vtMBs\endcsname{18.8}
\expandafter\def\csname vmMC\endcsname{+2.520}
\expandafter\def\csname vmMClo\endcsname{+2.311}
\expandafter\def\csname vmMChi\endcsname{+2.729}
\expandafter\def\csname vtMCg\endcsname{284.2}
\expandafter\def\csname vtMCs\endcsname{129.0}
\expandafter\def\csname vk8vs64\endcsname{2.29}
\expandafter\def\csname vk82\endcsname{+0.098}
\expandafter\def\csname vk82p\endcsname{0.409}
\expandafter\def\csname vkpar\endcsname{3.9}
\expandafter\def\csname vktime\endcsname{2.5}
\expandafter\def\csname vpredP\endcsname{1.0\times10^{-5}}
\expandafter\def\csname vtin64r\endcsname{97.7}
\expandafter\def\csname vanA\endcsname{+0.115}
\expandafter\def\csname vanAp\endcsname{0.4313}
\expandafter\def\csname vansA\endcsname{+0.647}
\expandafter\def\csname vansAT\endcsname{0.8259}
\expandafter\def\csname vanuA\endcsname{+0.870}
\expandafter\def\csname vanuAp\endcsname{0.0044}
\expandafter\def\csname vanuAT\endcsname{0.9545}
\expandafter\def\csname vhcA\endcsname{0.03983}
\expandafter\def\csname vhaA\endcsname{0.00282}
\expandafter\def\csname vflA\endcsname{14}
\expandafter\def\csname vlaA\endcsname{0.1540}
\expandafter\def\csname vlcA\endcsname{0.0638}
\expandafter\def\csname vbcA\endcsname{+0.755}
\expandafter\def\csname vanB\endcsname{-0.682}
\expandafter\def\csname vanBp\endcsname{0.0017}
\expandafter\def\csname vansB\endcsname{+0.472}
\expandafter\def\csname vansBT\endcsname{0.4319}
\expandafter\def\csname vanuB\endcsname{+0.297}
\expandafter\def\csname vanuBp\endcsname{0.0810}
\expandafter\def\csname vanuBT\endcsname{0.0976}
\expandafter\def\csname vhcB\endcsname{0.02132}
\expandafter\def\csname vhaB\endcsname{0.00073}
\expandafter\def\csname vflB\endcsname{29}
\expandafter\def\csname vlaB\endcsname{0.0173}
\expandafter\def\csname vlcB\endcsname{0.0132}
\expandafter\def\csname vbcB\endcsname{+0.978}
\expandafter\def\csname vanC\endcsname{-0.048}
\expandafter\def\csname vanCp\endcsname{0.7449}
\expandafter\def\csname vansC\endcsname{+0.032}
\expandafter\def\csname vansCT\endcsname{0.0025}
\expandafter\def\csname vanuC\endcsname{-0.017}
\expandafter\def\csname vanuCp\endcsname{0.8369}
\expandafter\def\csname vanuCT\endcsname{0.0007}
\expandafter\def\csname vhcC\endcsname{0.01147}
\expandafter\def\csname vhaC\endcsname{0.00018}
\expandafter\def\csname vflC\endcsname{63}
\expandafter\def\csname vlaC\endcsname{0.0080}
\expandafter\def\csname vlcC\endcsname{0.0074}
\expandafter\def\csname vbcC\endcsname{+0.032}
\expandafter\def\csname vconstT64\endcsname{0.0188}
\expandafter\def\csname vspanlo\endcsname{0.228}
\expandafter\def\csname vspanhi\endcsname{0.272}
\expandafter\def\csname vspanpct\endcsname{9}

\title{\vspace{-1.3cm}\textbf{Where Does the Gain in Dynamic Convolution Come From?}\\[3pt]
\large A controlled study of routing, capacity and width on CIFAR-100 and Tiny-ImageNet}
\author{Ankit Kumar\\\small Mechatronics Department, University of Agder, Grimstad, Norway}
\date{\small 24 September 2026}

\begin{document}
\maketitle
\vspace{-0.6cm}

\begin{abstract}
\noindent
Dynamic convolution replaces one kernel with a bank of $K$ kernels mixed
by input-dependent weights. Its gains are usually credited to
input-adaptive kernel selection, but published comparisons set it
against a network with $K$ times fewer parameters, so routing and
capacity are confounded. We separate them with a constant-mixture
control, a mismatched-input control and a frozen random router, across
\V{nruns} training runs on CIFAR-100 and Tiny-ImageNet. Four results hold
throughout. The kernel bank on its own contributes nothing: in all five
conditions tested, a network whose $K$ kernels are mixed uniformly is
statistically equivalent to a plain network at a $\pm$\SI{0.5}{pp}
margin. Where dynamic convolution helps, the routing term accounts for
the effect. At matched parameter budgets, a plain wider network beats the
routed one at all three budgets we could match, by 2.3 to
\SI{4.3}{pp}, at a fraction of the training cost. And routing entropy,
the statistic used to show that a router is active, does not track the
benefit: a router \V{flA} times flatter by entropy delivered a larger gain
than a sharper one. The size of the routing effect is less stable. On
CIFAR-100 it vanishes at the standard width of 64, under both a constant
and an annealed softmax temperature. On Tiny-ImageNet it does not
(\V{effTB}\,pp at width 64). Before running that experiment we recorded
a prediction that the effect would vanish there because the network
memorises its training data. The prediction failed, and we report it as
failed.
\end{abstract}

\section{Introduction}

A dynamic convolution layer stores $K$ kernels and combines them per
input. For an input feature map $F$ and kernel bank
$\{\mathcal{W}_k\}_{k=1}^{K}$, the layer computes
\begin{equation}
\mathcal{W}_{\text{dyn}}(F)=\sum_{k=1}^{K} z_k(F)\,\mathcal{W}_k,
\qquad
\zv=\operatorname{softmax}\!\left(\frac{g(\phi(F))}{\tau}\right),
\label{eq:dynconv}
\end{equation}
where $\phi$ summarises the feature map, $g$ is a small multilayer
perceptron and $\tau$ is a softmax temperature. Nearly every published
variant sets $\phi$ to global average pooling
\citep{yang2019condconv,chen2020dynamic}.

The usual account of why this helps is that $\zv$ adapts the kernel to
the input. A second account is available and is rarely tested. The bank
multiplies the stored convolution parameters by $K$, and published
comparisons are made against a network with $K$ times fewer parameters.
Under that comparison, a gain from routing and a gain from extra
capacity look identical.

This report separates the two. We compare the standard routed layer
against a constant mixture that has the same bank and no routing, a
router fed a different image's descriptor, a router frozen at random
initialisation, and a plain network. The comparison spans three widths,
three bank sizes, two temperature schedules and two datasets, for
\V{nruns} training runs and \V{gpuh} GPU-hours on one RTX 4070 Ti. Every
comparison is paired by seed and tested for equivalence as well as for
difference, so that a null result is a claim with a stated margin
rather than a failure to reject.

Four findings hold across everything we ran:
\begin{enumerate}\setlength\itemsep{1pt}
\item The kernel bank contributes nothing on its own, in all five
conditions tested across both datasets (Section~\ref{sec:decomp}).
\item Where dynamic convolution helps, the routing term accounts for the
effect.
\item At matched parameter budgets, spending the parameters on width
beats spending them on a kernel bank, at all three budgets we could
match (Section~\ref{sec:matched}).
\item Routing entropy does not measure whether routing is useful
(Section~\ref{sec:entropy}).
\end{enumerate}

The size of the routing effect is less stable than any of these. On
CIFAR-100 it peaks at width 32 and vanishes at width 64. We attributed
that null to the network memorising its training set, and before
extending the study to Tiny-ImageNet at width 64 we wrote down a
prediction that the effect would vanish there as well. It did not. We
report that prediction as failed (Section~\ref{sec:tin}) and treat the
CIFAR-100 width-64 null as an observation we cannot yet explain.

\section{Related work}

\paragraph{Dynamic convolution.} CondConv \citep{yang2019condconv} and
DY-Conv \citep{chen2020dynamic} introduced input-dependent mixtures of
kernels. ODConv \citep{li2022odconv} extends the attention across several
kernel dimensions, and KernelWarehouse \citep{li2024kernelwarehouse}
reorganises the bank so that far larger $K$ becomes affordable. The
reported gains are large: for a ResNet-18 on ImageNet at four times the
kernel parameters, \citet{li2024kernelwarehouse} report top-1
improvements of \SI{3.38}{pp} for DY-Conv, \SI{4.01}{pp} for ODConv and
\SI{5.61}{pp} for KernelWarehouse under their training recipe. FDConv
\citep{chen2025fdconv} observes that the parallel kernels of DY-Conv tend
toward highly similar frequency responses, so that the extra parameters
buy little adaptability. Our finding that the bank contributes nothing on
its own is consistent with that observation. Spatially resolved variants
such as DRConv \citep{chen2021drconv} route on local rather than global
signals and fall outside the controls used here.

\paragraph{Optimisation difficulty.} \citet{li2021dcd} identify the joint
optimisation of dynamic attention and static kernels as the central
difficulty of dynamic convolution and propose a matrix decomposition that
eases it. DY-Conv addresses the same difficulty by annealing the softmax
temperature from 30 to 1 early in training \citep{chen2020dynamic}.
Section~\ref{sec:temp} tests that schedule.

\paragraph{Re-parameterisation.} A constant mixture of $K$ kernels is a
linear re-parameterisation of a single kernel. Such re-parameterisations
can change optimisation without changing the function class
\citep{arora2018optimization}, and ACNet and RepVGG exploit this by
training a richer block and merging it into one kernel at inference
\citep{ding2019acnet,ding2021repvgg}. Our \texttt{uniform} arm is exactly
such a mergeable block.

\paragraph{Memorisation.} Networks of the size used here can fit their
training sets completely, even under random labels \citep{zhang2017}.
That motivated the saturation hypothesis tested in Section~\ref{sec:tin}.

\paragraph{Equivalence testing.} The two one-sided tests procedure
\citep{schuirmann1987} asserts that a difference lies within a stated
margin; \citet{lakens2017} give a practical treatment. It is standard in
bioequivalence studies and rare in this literature, where a large
$p$-value from a difference test is often read as evidence of no effect.

\section{Method}\label{sec:method}

\subsection{Routing signals}

All arms share Eq.~\eqref{eq:dynconv}, the backbone, the optimiser and the
seeds, and differ only in how $\zv$ is produced. Write
$\mu(F)\in\mathbb{R}^{C}$ for the channel-wise global average.

\begin{description}\setlength\itemsep{2pt}
\item[\texttt{static}] No dynamic layer: a plain ResNet-18, $K=1$.
\item[\texttt{uniform}] $g\equiv\mathbf{0}$, so $\zv=(1/K,\dots,1/K)$ for
every input. The bank exists and is trained; nothing is routed.
\item[\texttt{gap}] $\zv=\operatorname{softmax}(g(\mu(F))/\tau)$ with $g$
trained. The standard layer.
\item[\texttt{shuffled}] As \texttt{gap}, but sample $i$ is routed with
the descriptor of sample $\pi(i)$ for a permutation $\pi$ redrawn on every
forward pass, at evaluation as well as training. The distribution of
$\zv$ is preserved and the pairing between image and kernel is
destroyed.
\item[\texttt{frozen\_random}] As \texttt{gap}, with every weight of $g$
frozen at its random initialisation. Routing varies with the input but
is never optimised.
\end{description}

\subsection{The constant mixture is a single kernel}

\begin{proposition}\label{prop:collapse}
If $\zv=(1/K,\dots,1/K)$ for every input, the layer computes
$F*\bar{\mathcal{W}}$ with $\bar{\mathcal{W}}=\tfrac1K\sum_k\mathcal{W}_k$,
and its effective parameter count equals that of one static kernel.
\end{proposition}
\begin{proof}
The mixture in Eq.~\eqref{eq:dynconv} is constant, so
$\mathcal{W}_{\text{dyn}}=\bar{\mathcal{W}}$ for every $F$.
\end{proof}

At inference \texttt{uniform} is therefore a plain ResNet-18. At width 64
it stores 44.18M parameters and has an effective count of 11.22M. It
spans the same function class as \texttt{static}, so any difference
between the two is an effect of optimisation. That makes it the right
control for the extra parameters of the bank.

\subsection{Two exact implementations}

Convolution is linear in its weights, so
\begin{equation}
F*\Big(\sum_k z_k\mathcal{W}_k\Big)=\sum_k z_k\,(F*\mathcal{W}_k).
\label{eq:lin}
\end{equation}
The left side builds a per-sample kernel and runs one grouped
convolution; at width 64 and batch 128 its aggregated weights alone take
about \SI{1.2}{\giga\byte} per stage-4 layer. The right side runs $K$
ordinary convolutions and mixes the outputs, which costs $K$ times the
arithmetic but no per-sample weights. Every run uses the right side. The
test suite asserts that the two agree to $10^{-10}$ in double precision
across four routing signals.

\subsection{Training protocol}

\paragraph{Data.} CIFAR-100 \citep{krizhevsky2009}, $32\times32$, 100
classes, at widths 16, 32 and 64. Tiny-ImageNet \citep{le2015tiny},
$64\times64$, 200 classes, at widths 32 and 64. The stem is the same on
both: a $3\times3$ stride-1 convolution with no max-pool. At
$64\times64$ this leaves an $8\times8$ final feature map rather than
$4\times4$ and roughly quadruples the cost of the deeper stages, but
changing the stem would change the network between datasets. Crop
padding is 4 pixels at 32 and 8 at 64.

\paragraph{Network and recipe.} ResNet-18 \citep{he2016deep} in which all
16 $3\times3$ convolutions inside residual blocks are dynamic; the stem
and the $1\times1$ shortcuts stay static. $K=4$ unless stated. SGD with
Nesterov momentum 0.9, learning rate 0.1, weight decay $5\times10^{-4}$,
batch 128, five warm-up epochs then cosine decay over 150 epochs.
Augmentation is random cropping and horizontal flipping only. Mixed
precision in bfloat16. Six seeds, 0 to 5, shared by every arm. Batch 128
held at Tiny-ImageNet width 64 after a one-epoch memory probe, so the two
Tiny-ImageNet widths are directly comparable.

\paragraph{Temperature.} $\tau=1$ throughout, except in
Section~\ref{sec:temp}, where $\tau$ falls linearly from 30 to 1 over the
first ten epochs and stays at 1 thereafter.

\section{Statistical protocol}\label{sec:stats}

\paragraph{Equivalence.} Several conclusions are null conclusions, and a
large $p$ from a difference test does not establish one. For paired
differences $d_i$ with mean $\bar d$ and standard error $s_{\bar d}$ we
use two one-sided tests against a margin $\delta$,
\begin{equation}
p_{\text{TOST}}=\max\Big\{1-T_{n-1}\!\Big(\tfrac{\bar d+\delta}{s_{\bar d}}\Big),\;
T_{n-1}\!\Big(\tfrac{\bar d-\delta}{s_{\bar d}}\Big)\Big\},
\end{equation}
and read $p_{\text{TOST}}<0.05$ as a positive claim that the true
difference lies within $\pm\delta$. We fixed $\delta=\SI{0.5}{pp}$ before
the first run. The pre-registration justified this by reference to
reported gains at CIFAR scale without citing a source. The nearest
documented figures we have since found are the ImageNet ResNet-18 gains
of 3.38 to \SI{5.61}{pp} quoted above, and the margin sits an
order of magnitude below them.

\paragraph{Why six seeds.} The two-sided Wilcoxon signed-rank test
\citep{wilcoxon1945} on $n$ pairs admits only $2^n$ equally likely sign
assignments, so its smallest attainable $p$ is $2^{1-n}$: 0.25 at $n=3$,
0.0625 at $n=5$, 0.031 at $n=6$. Six is the smallest count at which it
can reach 0.05. We report it beside the paired $t$-test throughout.

\paragraph{Power.} For $\texttt{gap}-\texttt{uniform}$ on CIFAR-100 at
width 64 the paired standard deviation is \V{powsd}\,pp, giving a
minimum detectable difference of \V{mde}\,pp at $n=6$. Simulated
power ($2\times10^5$ draws) is \V{pwA}\,\% for a true effect of
\SI{0.25}{pp}, \V{pwB}\,\% at \SI{0.5}{pp} and
\V{pwC}\,\% at \SI{1.0}{pp}. The width-64 null on CIFAR-100
is not a power artefact.

\paragraph{Pre-registration.} Arms, recipe, primary metric, seeds, margin
and the decision table were written down before the first run finished.
The primary metric is final-epoch top-1 accuracy, since best-epoch
accuracy maximises over a noisy sequence. The Tiny-ImageNet width-64
prediction of Section~\ref{sec:tin} was recorded in writing, dated,
before any of those runs. An expectation about Tiny-ImageNet width 32 was
stated before those runs but not recorded; we give it correspondingly
less weight.

\section{CIFAR-100 at the standard width}\label{sec:w64}

\begin{figure}[t]\centering
\includegraphics[width=\textwidth]{fig/f3_w64.pdf}
\caption{CIFAR-100, width 64. (a) Every run of the five arms; bars are
arm means and the shaded band is $\pm$\SI{0.5}{pp} around
\texttt{static}. (b) Router flatness $1-\tilde H$ per seed, log scale;
\texttt{uniform} is exactly zero and is drawn at the axis floor.}
\label{fig:w64}
\end{figure}

At width 64 every arm is equivalent to a plain ResNet-18
(Table~\ref{tab:vsstatic}, Figure~\ref{fig:w64}). The largest gap between
arm means is \V{spread64}\,pp, against a median within-arm seed
standard deviation of \V{seedsd64}\,pp. A one-way ANOVA across the
five arms gives $F(4,25)=\V{anovaF}$, $p=\V{anovaP}$.

A router whose weights were never trained matches a trained one
($\texttt{frozen\_random}-\texttt{gap}=\V{frzVgap}$\,pp, $p=\V{frzVgapP}$,
$p_{\text{TOST}}=\V{frzVgapT}$), and so does a router that routes each
image with another image's descriptor. The trained router is not
degenerate: its per-seed mean entropy stays below that of the controls
(Figure~\ref{fig:w64}b). It discriminates between kernels, and at this
width the discrimination changes nothing.

\section{The capacity sweep on CIFAR-100}\label{sec:capacity}

\begin{figure}[t]\centering
\includegraphics[width=\textwidth]{fig/f1_sweep.pdf}
\caption{CIFAR-100. (a) Accuracy against base width, $\pm1$ standard
deviation over six seeds. (b) Paired difference from \texttt{static} with
\SI{95}{\percent} confidence intervals.}
\label{fig:sweep}
\end{figure}

Reducing width changes the answer (Figure~\ref{fig:sweep}).
$\texttt{gap}-\texttt{static}$ is \V{effCA}\,pp at width 16
($p\V{effCAp}$), \V{effCB}\,pp at width 32 ($p\V{effCBp}$) and
\V{effCC}\,pp at width 64 ($p\V{effCCp}$). The effect is not monotone in
capacity: it peaks at width 32. \texttt{uniform} is equivalent to
\texttt{static} at all three widths. Figure~\ref{fig:allruns} in the
appendix shows every run.

\section{The kernel bank contributes nothing}\label{sec:decomp}

\begin{figure}[t]\centering
\includegraphics[width=0.84\textwidth]{fig/f4_decomp.pdf}
\caption{Decomposition of the effect of dynamic convolution into the
bank ($\texttt{uniform}-\texttt{static}$) and routing
($\texttt{gap}-\texttt{uniform}$), five conditions, \SI{95}{\percent}
CIs.}
\label{fig:decomp}
\end{figure}

Splitting the total effect $\texttt{gap}-\texttt{static}$ into a bank
term and a routing term separates the two accounts of the introduction.
Table~\ref{tab:decomp} and Figure~\ref{fig:decomp} give the split in all
five conditions that have a \texttt{uniform} arm.

The bank term is equivalent to zero in every one of them, on both
datasets, with $p_{\text{TOST}}$ between 0.014 and 0.025. The largest
bank term observed is \V{bankCB}\,pp. Since the bank contributes nothing
within the margin, the total effect is the routing effect: on
Tiny-ImageNet at width 64, for example, routing is \V{routTB}\,pp of a
\V{effTB}\,pp total. At CIFAR-100 width 16 the bank term is slightly
negative (\V{bankCA}\,pp, equivalent) and the routing term therefore
exceeds the total, which is why we report no percentage share.

\begin{table}[t]\centering\small
\caption{Bank and routing terms, paired by seed. The bank term is
equivalent to zero in all five conditions.}
\label{tab:decomp}
\begin{tabular}{l r rll rl rl}
\toprule
& & \multicolumn{3}{c}{bank (\texttt{uniform}$-$\texttt{static})}
& \multicolumn{2}{c}{routing (\texttt{gap}$-$\texttt{uniform})}
& \multicolumn{2}{c}{total} \\
\cmidrule(lr){3-5}\cmidrule(lr){6-7}\cmidrule(lr){8-9}
dataset & width & $\Delta$ (pp) & $p_{\text{TOST}}$ & verdict & $\Delta$ (pp) & $p$ & $\Delta$ (pp) & $p$ \\
\midrule
CIFAR-100 & 16 & -0.223 & 0.0251 & equivalent & +0.755 & 0.0023 & +0.532 & 0.0023 \\
CIFAR-100 & 32 & +0.175 & 0.0242 & equivalent & +0.978 & 0.0010 & +1.153 & $<10^{-4}$ \\
CIFAR-100 & 64 & +0.048 & 0.0142 & equivalent & +0.032 & 0.7752 & +0.080 & 0.6158 \\
Tiny-ImageNet & 32 & +0.107 & 0.0218 & equivalent & +2.078 & $<10^{-4}$ & +2.185 & $<10^{-4}$ \\
Tiny-ImageNet & 64 & +0.053 & 0.0157 & equivalent & +2.228 & $<10^{-4}$ & +2.282 & $<10^{-4}$ \\
\bottomrule
\end{tabular}
\end{table}

\section{Bank size}\label{sec:ksweep}

\begin{figure}[t]\centering
\includegraphics[width=0.84\textwidth]{fig/f6_k.pdf}
\caption{CIFAR-100, width 32. (a) Effect over \texttt{static} for
$K\in\{2,4,8\}$, \SI{95}{\percent} CIs. (b) Stored parameters and
training time.}
\label{fig:k}
\end{figure}

The width-32 peak was measured at $K=4$; Table~\ref{tab:k} repeats it at
$K=2$ and $K=8$. All three bank sizes beat \texttt{static} by a similar
margin, and no pair of them is distinguishable (Table~\ref{tab:kpairs}):
$K=8$ over $K=2$ is \V{k82}\,pp ($p=\V{k82p}$). The benefit is already
complete at $K=2$. Going from $K=2$ to $K=8$ multiplies stored parameters
by \V{kpar} and training time by \V{ktime} (Figure~\ref{fig:k}).

\begin{table}[t]\centering\small
\caption{Bank size at CIFAR-100 width 32. Differences are against the
\texttt{static} row.}
\label{tab:k}
\begin{tabular}{l rr rl l rr}
\toprule
$K$ & accuracy & sd & $\Delta$ (pp) & \SI{95}{\percent} CI & $p$ & params (M) & min/run \\
\midrule
2 & 75.950 & 0.148 & +1.078 & [+0.758, +1.399] & 0.0003 & 5.64 & 27.9 \\
4 & 76.025 & 0.156 & +1.153 & [+0.940, +1.367] & $<10^{-4}$ & 11.13 & 39.2 \\
8 & 76.048 & 0.242 & +1.177 & [+0.749, +1.604] & 0.0009 & 22.12 & 69.3 \\
\midrule \texttt{static} & 74.872 & 0.262 & -- & -- & -- & 2.82 & 8.6 \\
\bottomrule
\end{tabular}
\end{table}

\begin{table}[t]\centering\small
\caption{Pairwise comparisons between bank sizes.}
\label{tab:kpairs}
\begin{tabular}{l rl ll l}
\toprule
comparison & $\Delta$ (pp) & \SI{95}{\percent} CI & $p$ & $p_{\text{TOST}}$ & verdict \\
\midrule
$K=4$ vs $K=2$ & +0.075 & [-0.087, +0.237] & 0.2878 & 0.0005 & equivalent \\
$K=8$ vs $K=4$ & +0.023 & [-0.286, +0.332] & 0.8538 & 0.0053 & equivalent \\
$K=8$ vs $K=2$ & +0.098 & [-0.182, +0.379] & 0.4091 & 0.0072 & equivalent \\
\bottomrule
\end{tabular}
\end{table}

\section{Tiny-ImageNet, and a failed prediction}\label{sec:tin}

\begin{figure}[t]\centering
\includegraphics[width=\textwidth]{fig/f7_tin.pdf}
\caption{(a) All 36 Tiny-ImageNet runs. (b) $\texttt{gap}-\texttt{static}$
by width on both datasets, \SI{95}{\percent} CIs. The CIFAR-100
width-64 null does not appear on Tiny-ImageNet.}
\label{fig:tin}
\end{figure}

\subsection{Width 32}

On Tiny-ImageNet at width 32, $\texttt{gap}-\texttt{static}$ is
\V{effTA}\,pp (\SI{95}{\percent} CI [\V{effTAlo}, \V{effTAhi}],
$p\V{effTAp}$), against \V{effCB}\,pp on CIFAR-100 at the same width. The
bank term is \V{bankTA}\,pp and equivalent to zero. Before these runs we
expected the effect to be larger than on CIFAR-100, since the static arm
trains to a higher loss (\V{lossTA} against \V{lossCB}) and so sits
further from memorisation. That expectation held, but it was stated and
not recorded, it predicted only a direction, and we do not treat it as
strong evidence.

\subsection{Width 64: the recorded prediction}

The saturation hypothesis says the routing effect vanishes once a network
memorises its training data. On CIFAR-100 the width-64 null coincides
with a final training loss of \V{lossCC}. Before running Tiny-ImageNet at
width 64 we recorded, in writing and dated, the prediction in
Table~\ref{tab:predict}: that the effect would fall below \SI{0.5}{pp}
and be equivalent to zero, and that the static arm would train to a loss
below 0.05.

\begin{table}[t]\centering\small
\caption{The prediction recorded before the Tiny-ImageNet width-64 runs,
and the outcome. Both conditions were stated as required.}
\label{tab:predict}
\begin{tabular}{l p{5.4cm} p{4.6cm} l}
\toprule
& prediction & measured & outcome \\
\midrule
(a) & \texttt{gap} $-$ \texttt{static} below \SI{0.5}{pp}, equivalent by TOST & +2.282 pp, $p=1.0\times10^{-5}$, $p_{\text{TOST}}=1.0000$ & failed \\
(b) & static final training loss below 0.05 & 0.021 & held \\
\bottomrule
\end{tabular}
\end{table}

The loss condition held (\V{lossTB}) and the effect did not vanish:
$\texttt{gap}-\texttt{static}=\V{effTB}$\,pp, \SI{95}{\percent} CI
[\V{effTBlo}, \V{effTBhi}], $p=\V{predP}$. The prediction failed. The
bank term is again equivalent to zero (\V{bankTB}\,pp,
$p_{\text{TOST}}=\V{bankTBT}$), so the effect is carried by routing, and
it is as large at width 64 as at width 32.

\subsection{What remains of the saturation account}

\begin{figure}[t]\centering
\includegraphics[width=\textwidth]{fig/f8_prediction.pdf}
\caption{(a) Training loss of static and gap arms, seed 0, log scale.
(b) Routing effect against the final training loss of the static arm in
all five conditions. The shaded band is the region the prediction
associated with a null.}
\label{fig:pred}
\end{figure}

Table~\ref{tab:satur} orders the five conditions by the final training
loss of the static arm. The effects do not follow that order
(Figure~\ref{fig:pred}b). Tiny-ImageNet at width 64 trains to a loss
close to CIFAR-100 at width 32 and shows twice the effect; CIFAR-100 at
width 16 trains to the highest loss and shows one of the smallest. A
threshold near 0.01 would separate the one null from the four positive
results, but it would be chosen after seeing the data, and we do not
adopt it.

What survives is narrower. On CIFAR-100 the routing effect vanishes at
width 64, and that result is well powered and holds under two temperature
schedules. It does not replicate on Tiny-ImageNet, and training loss does
not explain it. Plausible differences between the datasets include the
input resolution, the final feature-map size and the number of classes;
none of them was isolated here.

\begin{table}[t]\centering\small
\caption{Five conditions ordered by the final training loss of the
static arm.}
\label{tab:satur}
\begin{tabular}{l r r rl l}
\toprule
dataset & width & static train loss & $\texttt{gap}-\texttt{static}$ (pp) & \SI{95}{\percent} CI & verdict \\
\midrule
CIFAR-100 & 64 & 0.0080 & +0.080 & [-0.305, +0.465] & equivalent \\
CIFAR-100 & 32 & 0.0170 & +1.153 & [+0.940, +1.367] & higher \\
Tiny-ImageNet & 64 & 0.0210 & +2.282 & [+1.952, +2.611] & higher \\
Tiny-ImageNet & 32 & 0.1194 & +2.185 & [+1.781, +2.589] & higher \\
CIFAR-100 & 16 & 0.1643 & +0.532 & [+0.292, +0.771] & higher \\
\bottomrule
\end{tabular}
\end{table}

\section{Matched parameter budgets}\label{sec:matched}

\begin{figure}[t]\centering
\includegraphics[width=\textwidth]{fig/f5_matched.pdf}
\caption{Accuracy against stored parameters. Dotted arrows join routed
and plain networks matched to within \SI{1}{\percent} on parameter
count.}
\label{fig:matched}
\end{figure}

Because widths differ by a factor of two and the bank multiplies
convolution parameters by four, a routed network at one width stores
almost exactly as many parameters as a plain network at the next width
up. That gives three matched comparisons without any extra runs
(Table~\ref{tab:matched}, Figure~\ref{fig:matched}).

The plain network wins all three: by \V{mMA}\,pp at 2.8M
parameters and \V{mMB}\,pp at 11.2M on CIFAR-100, and by
\V{mMC}\,pp at 11.2M on Tiny-ImageNet. It also trains faster in
each case; on Tiny-ImageNet, \V{tMCs}\,min per run against
\V{tMCg}\,min. A fourth, unmatched comparison points the same
way: the $K=8$ network at width 32 stores 22.1M parameters, twice the
plain width-64 network, and trails it by \V{k8vs64}\,pp.

\begin{table}[t]\centering\small
\caption{Routed against plain networks at matched parameter counts,
paired by seed. Minutes are routed / plain per run.}
\label{tab:matched}
\begin{tabular}{l l l rr rl l r}
\toprule
dataset & routed & plain & routed & plain & $\Delta$ (pp) & \SI{95}{\percent} CI & $p$ & min/run \\
\midrule
CIFAR-100 & \texttt{gap} w16 (2.79M) & \texttt{static} w32 (2.82M) & 70.527 & 74.872 & +4.345 & [+4.103, +4.587] & $<10^{-4}$ & 33.7 / 8.6 \\
CIFAR-100 & \texttt{gap} w32 (11.13M) & \texttt{static} w64 (11.22M) & 76.025 & 78.338 & +2.313 & [+1.870, +2.757] & $<10^{-4}$ & 39.2 / 18.8 \\
Tiny-ImageNet & \texttt{gap} w32 (11.16M) & \texttt{static} w64 (11.27M) & 63.745 & 66.265 & +2.520 & [+2.311, +2.729] & $<10^{-4}$ & 284.2 / 129.0 \\
\bottomrule
\end{tabular}
\end{table}

\section{Temperature schedule}\label{sec:temp}

\begin{figure}[t]\centering
\includegraphics[width=\textwidth]{fig/f9_temp.pdf}
\caption{CIFAR-100. (a) Routing effect under a constant and an annealed
temperature, \SI{95}{\percent} CIs. (b) Router flatness. (c) Flatness
against benefit over \texttt{uniform}; the two schedules do not lie on
one curve.}
\label{fig:temp}
\end{figure}

With $\tau=1$ the softmax can be close to one-hot from the first step,
starving most of the bank of gradient. DY-Conv anneals $\tau$ from 30 to
1 for this reason \citep{chen2020dynamic}. If the CIFAR-100 width-64
null came from an undertrained bank, it would say nothing about dynamic
convolution, so we repeated \texttt{gap} at all three widths with the
annealed schedule (Table~\ref{tab:temp}).

The width-64 null survives: annealed \texttt{gap} against \texttt{static}
is \V{ansC}\,pp with $p_{\text{TOST}}=\V{ansCT}$, a tighter equivalence
than the \V{constT64} obtained at $\tau=1$. Annealing is otherwise
neutral at width 64 (\V{anC}\,pp, $p=\V{anCp}$) and width 16
(\V{anA}\,pp, $p=\V{anAp}$), and harmful at width 32 (\V{anB}\,pp,
$p=\V{anBp}$), where it also leaves a higher final training loss
(\V{laB} against \V{lcB}). Backbone, dataset, bank size and schedule
length all differ from the DY-Conv setting, so this is not evidence
against their result; it is evidence that the CIFAR-100 null does not
depend on the choice of schedule.

At width 64 the annealed router went flat during the high-$\tau$ phase
and stayed flat for the remaining 140 epochs: it is \V{flC} times flatter
than the constant router ($1-\tilde H=\V{haC}$) and equivalent to
\texttt{uniform} (\V{anuC}\,pp, $p_{\text{TOST}}=\V{anuCT}$). This follows
from Proposition~\ref{prop:collapse}. A flat mixture trains the bank
toward its own mean, and a homogeneous bank leaves the router nothing to
discriminate between. At width 16 the same schedule did not produce that
outcome, which the next section takes up.

\begin{table}[t]\centering\small
\caption{Temperature ablation on CIFAR-100, \texttt{gap} arm. The last
two columns test whether the annealed router behaves like the constant
mixture.}
\label{tab:temp}
\begin{tabular}{r rr rl r rl rl}
\toprule
& \multicolumn{2}{c}{accuracy} & \multicolumn{2}{c}{annealed $-$ constant}
& \multicolumn{3}{c}{vs \texttt{static}} & \multicolumn{2}{c}{annealed vs \texttt{uniform}} \\
\cmidrule(lr){2-3}\cmidrule(lr){4-5}\cmidrule(lr){6-8}\cmidrule(lr){9-10}
width & $\tau=1$ & annealed & $\Delta$ & $p$ & $\tau=1$ & annealed & $p_{\text{TOST}}$ & $\Delta$ & verdict \\
\midrule
16 & 70.527 & 70.642 & +0.115 & 0.4313 & +0.532 & +0.647 & 0.8259 & +0.870 & higher \\
32 & 76.025 & 75.343 & -0.682 & 0.0017 & +1.153 & +0.472 & 0.4319 & +0.297 & inconclusive \\
64 & 78.418 & 78.370 & -0.048 & 0.7449 & +0.080 & +0.032 & 0.0025 & -0.017 & equivalent \\
\bottomrule
\end{tabular}
\end{table}

\section{Routing entropy does not measure routing usefulness}\label{sec:entropy}

\begin{figure}[t]\centering
\includegraphics[width=0.74\textwidth]{fig/f10_entropy.pdf}
\caption{CIFAR-100, width 16. The annealed router is \V{flA} times flatter
by entropy and delivers the larger benefit over \texttt{uniform}.}
\label{fig:ent}
\end{figure}

Normalised routing entropy $\tilde H=H(\zv)/\log K$ is the statistic
this literature reports to show that a router is doing something.
Table~\ref{tab:entben} sets it against the benefit each router delivers
over \texttt{uniform} in six conditions.

At width 16 the two schedules give a counterexample
(Figure~\ref{fig:ent}). The constant router has $1-\tilde H=\V{hcA}$ and
beats \texttt{uniform} by \V{bcA}\,pp. The annealed router has
$1-\tilde H=\V{haA}$, \V{flA} times flatter, and beats \texttt{uniform} by
\V{anuA}\,pp ($p=\V{anuAp}$). A router that looks nearly inactive on the
standard diagnostic delivered more than one that looks clearly active.

The reason is that entropy compresses the signal quadratically. Writing
$z_k=1/K+\varepsilon_k$ with $\sum_k\varepsilon_k=0$ and expanding to
second order,
\begin{equation}
1-\tilde H\;\approx\;\frac{K}{2\log K}\sum_k\varepsilon_k^{2},
\label{eq:ent}
\end{equation}
so the deviations that carry input-dependent information enter squared.
At $1-\tilde H=\V{haA}$ the root-mean-square deviation from $1/K$ is
about 0.022, so typical mixing weights lie in
[\V{spanlo}, \V{spanhi}], roughly \V{spanpct}\,\% of $1/K$.
Entropy reports that as 0.997.

Entropy is not meaningless. The annealed router at width 64, at
$1-\tilde H=\V{haC}$, has collapsed in the sense that matters: it is
equivalent to \texttt{uniform}. But across the six conditions of
Table~\ref{tab:entben}, sorting by $1-\tilde H$ does not sort the
benefit, and the two schedules do not lie on one curve
(Figure~\ref{fig:temp}c). A value of $\tilde H$ below one is not evidence
that routing is useful, and a value near one is not evidence that it is
not.

\begin{table}[t]\centering\small
\caption{Routing entropy against the benefit it is taken to indicate,
CIFAR-100. RMS dev is the root-mean-square deviation of the mixing
weights from $1/K$, recovered from Eq.~\eqref{eq:ent}.}
\label{tab:entben}
\begin{tabular}{r l rrr rl}
\toprule
width & $\tau$ & $\tilde H$ & $1-\tilde H$ & RMS dev & \texttt{gap}$-$\texttt{uniform} (pp) & $p$ \\
\midrule
16 & $1$ & 0.96017 & 0.03983 & 0.0831 & +0.755 & 0.0023 \\
16 & $30\!\to\!1$ & 0.99718 & 0.00282 & 0.0221 & +0.870 & 0.0044 \\
32 & $1$ & 0.97868 & 0.02132 & 0.0608 & +0.978 & 0.0010 \\
32 & $30\!\to\!1$ & 0.99927 & 0.00073 & 0.0113 & +0.297 & 0.0810 \\
64 & $1$ & 0.98853 & 0.01147 & 0.0446 & +0.032 & 0.7752 \\
64 & $30\!\to\!1$ & 0.99982 & 0.00018 & 0.0056 & -0.017 & 0.8369 \\
\bottomrule
\end{tabular}
\end{table}

\section{Threats to validity}

\paragraph{Scale.} Both datasets are low-resolution and every run fits on
one consumer GPU. Behaviour at ImageNet scale is untested and cannot be
tested on this hardware.

\paragraph{One backbone.} ResNet-18 throughout.

\paragraph{Not a reproduction of a published system.} The layer is a
dynamic convolution of the CondConv and DY-Conv form. No published
implementation of ODConv, FDConv or KernelWarehouse was run, and the
conclusions concern dynamic convolution of this form.

\paragraph{Global routing.} Every router here uses a global descriptor.
Spatially resolved routing is not covered.

\paragraph{One annealing schedule.} Section~\ref{sec:temp} compares a
constant temperature with one linear schedule.

\paragraph{Stored and effective parameters.} By
Proposition~\ref{prop:collapse} the \texttt{uniform} arm has the
effective capacity of a plain network. The matched-budget comparisons of
Section~\ref{sec:matched} use stored parameters, which is the count a
practitioner pays for.

\paragraph{The CIFAR-100 width-64 null is unexplained.} It is
well-powered and does not depend on the temperature schedule, and its
stated explanation failed a recorded test.

\paragraph{Provenance and corrections.} Final accuracies were transcribed
from the training logs, which print them verbatim, and every figure and
table is generated from those per-seed values. During the analysis, three
scripts grouped results by width alone, which merged runs from the two
datasets and from different bank sizes; this was detected when a
comparison paired one dataset's \texttt{uniform} arm with the other's
\texttt{static} arm, and every number in this report was recomputed from
per-seed values keyed on dataset, width, bank size and schedule. Summary
sentences printed by the analysis scripts were also found to over-read
non-significant orderings on four occasions; this report relies on the
per-comparison statistics only.

\section{Conclusion}

Dynamic convolution's gains, where they appear, come from routing and not
from the kernel bank: in all five conditions tested on two datasets, a
bank mixed uniformly is equivalent to a plain network. Two consequences
follow. Banks larger than $K=2$ buy nothing further, and at any
parameter budget we could match, the parameters are better spent on
width, by 2.3 to \SI{4.3}{pp} at a fraction of the training cost.

When routing helps is less settled. On CIFAR-100 it helps below width 64
and not at 64, under either temperature schedule. On Tiny-ImageNet it
helps at both widths tested. We proposed that memorisation of the
training set explains the CIFAR-100 null, recorded a prediction from that
account before testing it, and the prediction failed.

The diagnostic used to show that routing is active does not support that
use. A router \V{flA} times flatter by entropy delivered the larger
benefit, because entropy squares the deviations that carry the routing
signal.

\paragraph{Reproducibility.} \V{nruns} runs, \V{gpuh} GPU-hours, one RTX
4070 Ti. Code, the pre-registration, the dated prediction, the gate
ledger with recorded evidence, and every per-seed value accompany this
report.

\begin{thebibliography}{99}\small
\bibitem[Arora et~al.(2018)]{arora2018optimization} S.~Arora, N.~Cohen, and E.~Hazan. On the optimization of deep networks: Implicit acceleration by overparameterization. In \emph{ICML}, 2018.
\bibitem[Chen et~al.(2021)]{chen2021drconv} J.~Chen, X.~Wang, Z.~Guo, X.~Zhang, and J.~Sun. Dynamic region-aware convolution. In \emph{CVPR}, 2021.
\bibitem[Chen et~al.(2025)]{chen2025fdconv} L.~Chen, L.~Gu, L.~Li, C.~Yan, and Y.~Fu. Frequency dynamic convolution for dense image prediction. In \emph{CVPR}, pages 30178--30188, 2025.
\bibitem[Chen et~al.(2020)]{chen2020dynamic} Y.~Chen, X.~Dai, M.~Liu, D.~Chen, L.~Yuan, and Z.~Liu. Dynamic convolution: Attention over convolution kernels. In \emph{CVPR}, pages 11030--11039, 2020.
\bibitem[Ding et~al.(2019)]{ding2019acnet} X.~Ding, Y.~Guo, G.~Ding, and J.~Han. ACNet: Strengthening the kernel skeletons for powerful CNN via asymmetric convolution blocks. In \emph{ICCV}, 2019.
\bibitem[Ding et~al.(2021)]{ding2021repvgg} X.~Ding, X.~Zhang, N.~Ma, J.~Han, G.~Ding, and J.~Sun. RepVGG: Making VGG-style ConvNets great again. In \emph{CVPR}, 2021.
\bibitem[He et~al.(2016)]{he2016deep} K.~He, X.~Zhang, S.~Ren, and J.~Sun. Deep residual learning for image recognition. In \emph{CVPR}, 2016.
\bibitem[Krizhevsky(2009)]{krizhevsky2009} A.~Krizhevsky. Learning multiple layers of features from tiny images. Technical report, University of Toronto, 2009.
\bibitem[Lakens(2017)]{lakens2017} D.~Lakens. Equivalence tests: A practical primer for $t$ tests, correlations, and meta-analyses. \emph{Social Psychological and Personality Science}, 8(4):355--362, 2017.
\bibitem[Le and Yang(2015)]{le2015tiny} Y.~Le and X.~Yang. Tiny ImageNet visual recognition challenge. CS231N course report, Stanford University, 2015.
\bibitem[Li and Yao(2024)]{li2024kernelwarehouse} C.~Li and A.~Yao. KernelWarehouse: Rethinking the design of dynamic convolution. In \emph{ICML}, PMLR 235:29201--29221, 2024.
\bibitem[Li et~al.(2022)]{li2022odconv} C.~Li, A.~Zhou, and A.~Yao. Omni-dimensional dynamic convolution. In \emph{ICLR}, 2022.
\bibitem[Li et~al.(2021)]{li2021dcd} Y.~Li, Y.~Chen, X.~Dai, M.~Liu, D.~Chen, Y.~Yu, L.~Yuan, Z.~Liu, M.~Chen, and N.~Vasconcelos. Revisiting dynamic convolution via matrix decomposition. In \emph{ICLR}, 2021.
\bibitem[Schuirmann(1987)]{schuirmann1987} D.~J. Schuirmann. A comparison of the two one-sided tests procedure and the power approach for assessing the equivalence of average bioavailability. \emph{Journal of Pharmacokinetics and Biopharmaceutics}, 15(6):657--680, 1987.
\bibitem[Wilcoxon(1945)]{wilcoxon1945} F.~Wilcoxon. Individual comparisons by ranking methods. \emph{Biometrics Bulletin}, 1(6):80--83, 1945.
\bibitem[Yang et~al.(2019)]{yang2019condconv} B.~Yang, G.~Bender, Q.~V. Le, and J.~Ngiam. CondConv: Conditionally parameterized convolutions for efficient inference. In \emph{NeurIPS}, 2019.
\bibitem[Zhang et~al.(2017)]{zhang2017} C.~Zhang, S.~Bengio, M.~Hardt, B.~Recht, and O.~Vinyals. Understanding deep learning requires rethinking generalization. In \emph{ICLR}, 2017.
\end{thebibliography}

\clearpage
\appendix
\section{Complete results}

\begin{figure}[h]\centering
\includegraphics[width=\textwidth]{fig/f2_allruns.pdf}
\caption{Every CIFAR-100 run at $K=4$ and $\tau=1$; bars are arm means.}
\label{fig:allruns}
\end{figure}

\begin{table}[h]\centering\scriptsize
\caption{Every run at $K=4$ and $\tau=1$: final-epoch top-1 accuracy
(\si{\percent}) for each seed, with mean and sample standard deviation.}
\label{tab:allruns}
\begin{tabular}{l r l rrrrrr rr}
\toprule
dataset & width & arm & s0 & s1 & s2 & s3 & s4 & s5 & mean & sd \\
\midrule
CIFAR-100 & 16 & \texttt{static} & 69.91 & 69.75 & 69.86 & 70.46 & 70.00 & 69.99 & 69.995 & 0.246 \\
CIFAR-100 & 16 & \texttt{uniform} & 69.68 & 69.32 & 69.84 & 69.91 & 69.72 & 70.16 & 69.772 & 0.279 \\
CIFAR-100 & 16 & \texttt{gap} & 70.35 & 70.51 & 70.41 & 70.63 & 70.79 & 70.47 & 70.527 & 0.160 \\
CIFAR-100 & 32 & \texttt{static} & 74.81 & 74.65 & 74.53 & 74.92 & 75.11 & 75.21 & 74.872 & 0.262 \\
CIFAR-100 & 32 & \texttt{uniform} & 75.29 & 74.86 & 74.75 & 75.30 & 74.71 & 75.37 & 75.047 & 0.305 \\
CIFAR-100 & 32 & \texttt{gap} & 76.13 & 75.88 & 75.91 & 75.90 & 76.27 & 76.06 & 76.025 & 0.156 \\
CIFAR-100 & 64 & \texttt{static} & 78.05 & 78.67 & 78.63 & 78.27 & 78.00 & 78.41 & 78.338 & 0.284 \\
CIFAR-100 & 64 & \texttt{uniform} & 78.64 & 78.59 & 78.34 & 78.30 & 78.36 & 78.09 & 78.387 & 0.202 \\
CIFAR-100 & 64 & \texttt{gap} & 78.59 & 78.81 & 78.11 & 78.61 & 78.07 & 78.32 & 78.418 & 0.299 \\
CIFAR-100 & 64 & \texttt{shuffled} & 78.42 & 77.83 & 78.18 & 78.41 & 78.47 & 78.28 & 78.265 & 0.238 \\
CIFAR-100 & 64 & \texttt{frozen\_random} & 78.47 & 78.49 & 78.48 & 78.63 & 78.45 & 78.57 & 78.515 & 0.070 \\
Tiny-ImageNet & 32 & \texttt{static} & 61.93 & 61.61 & 61.57 & 61.46 & 61.27 & 61.52 & 61.560 & 0.217 \\
Tiny-ImageNet & 32 & \texttt{uniform} & 61.52 & 61.53 & 61.55 & 61.62 & 61.66 & 62.12 & 61.667 & 0.229 \\
Tiny-ImageNet & 32 & \texttt{gap} & 63.53 & 64.20 & 63.61 & 63.58 & 63.91 & 63.64 & 63.745 & 0.260 \\
Tiny-ImageNet & 64 & \texttt{static} & 66.29 & 66.59 & 65.88 & 65.97 & 66.63 & 66.23 & 66.265 & 0.309 \\
Tiny-ImageNet & 64 & \texttt{uniform} & 66.27 & 66.27 & 66.41 & 66.21 & 66.22 & 66.53 & 66.318 & 0.126 \\
Tiny-ImageNet & 64 & \texttt{gap} & 68.73 & 68.35 & 68.42 & 68.55 & 68.71 & 68.52 & 68.547 & 0.152 \\
\bottomrule
\end{tabular}
\end{table}

\begin{table}[h]\centering\small
\caption{Paired comparisons against \texttt{static}, $n=6$. $p_W$ is the
Wilcoxon signed-rank test, which cannot fall below 0.031 at this $n$.}
\label{tab:vsstatic}
\begin{tabular}{l r l r l ll r l l}
\toprule
dataset & width & arm & $\Delta$ (pp) & \SI{95}{\percent} CI & $p_t$ & $p_W$ & $d_z$ & $p_{\text{TOST}}$ & verdict \\
\midrule
CIFAR-100 & 16 & \texttt{uniform} & -0.223 & [-0.500, +0.054] & 0.0931 & 0.0938 & -0.85 & 0.0251 & equivalent \\
CIFAR-100 & 16 & \texttt{gap} & +0.532 & [+0.292, +0.771] & 0.0023 & 0.0312 & +2.33 & 0.6260 & higher \\
CIFAR-100 & 32 & \texttt{uniform} & +0.175 & [-0.146, +0.496] & 0.2206 & 0.3125 & +0.57 & 0.0242 & equivalent \\
CIFAR-100 & 32 & \texttt{gap} & +1.153 & [+0.940, +1.367] & $<10^{-4}$ & 0.0312 & +5.67 & 0.9997 & higher \\
CIFAR-100 & 64 & \texttt{uniform} & +0.048 & [-0.332, +0.429] & 0.7573 & 0.8438 & +0.13 & 0.0142 & equivalent \\
CIFAR-100 & 64 & \texttt{gap} & +0.080 & [-0.305, +0.465] & 0.6158 & 0.5625 & +0.22 & 0.0188 & equivalent \\
CIFAR-100 & 64 & \texttt{shuffled} & -0.073 & [-0.602, +0.456] & 0.7361 & 1.0000 & -0.15 & 0.0464 & equivalent \\
CIFAR-100 & 64 & \texttt{frozen\_random} & +0.177 & [-0.121, +0.474] & 0.1874 & 0.2188 & +0.62 & 0.0191 & equivalent \\
Tiny-ImageNet & 32 & \texttt{uniform} & +0.107 & [-0.270, +0.483] & 0.4994 & 0.6875 & +0.30 & 0.0218 & equivalent \\
Tiny-ImageNet & 32 & \texttt{gap} & +2.185 & [+1.781, +2.589] & $<10^{-4}$ & 0.0312 & +5.67 & 0.9999 & higher \\
Tiny-ImageNet & 64 & \texttt{uniform} & +0.053 & [-0.334, +0.441] & 0.7380 & 1.0000 & +0.14 & 0.0157 & equivalent \\
Tiny-ImageNet & 64 & \texttt{gap} & +2.282 & [+1.952, +2.611] & $<10^{-4}$ & 0.0312 & +7.26 & 1.0000 & higher \\
\bottomrule
\end{tabular}
\end{table}

\begin{table}[h]\centering\small
\caption{Configuration and cost. Router parameters are in thousands and
are included in the stored total. Training loss is the final-epoch mean
over six seeds.}
\label{tab:config}
\begin{tabular}{l r l rr r r}
\toprule
dataset & width & arm & params (M) & router (K) & train loss & min/run \\
\midrule
CIFAR-100 & 16 & \texttt{static} & 0.71 & 0.0 & 0.1643 & 9.6 \\
CIFAR-100 & 16 & \texttt{uniform} & 2.77 & 0.0 & 0.1653 & 11.4 \\
CIFAR-100 & 16 & \texttt{gap} & 2.79 & 21.3 & 0.0638 & 33.7 \\
CIFAR-100 & 32 & \texttt{static} & 2.82 & 0.0 & 0.0170 & 8.6 \\
CIFAR-100 & 32 & \texttt{uniform} & 11.06 & 0.0 & 0.0170 & 12.9 \\
CIFAR-100 & 32 & \texttt{gap} & 11.13 & 74.6 & 0.0132 & 39.2 \\
CIFAR-100 & 64 & \texttt{static} & 11.22 & 0.0 & 0.0080 & 18.8 \\
CIFAR-100 & 64 & \texttt{uniform} & 44.18 & 0.0 & 0.0081 & 25.2 \\
CIFAR-100 & 64 & \texttt{gap} & 44.47 & 288.0 & 0.0074 & 86.0 \\
CIFAR-100 & 64 & \texttt{shuffled} & 44.47 & 288.0 & 0.0083 & 92.3 \\
CIFAR-100 & 64 & \texttt{frozen\_random} & 44.18 & 0.0 & 0.0077 & 85.4 \\
Tiny-ImageNet & 32 & \texttt{static} & 2.85 & 0.0 & 0.1194 & 54.2 \\
Tiny-ImageNet & 32 & \texttt{uniform} & 11.09 & 0.0 & 0.1200 & 56.3 \\
Tiny-ImageNet & 32 & \texttt{gap} & 11.16 & 74.6 & 0.0420 & 284.2 \\
Tiny-ImageNet & 64 & \texttt{static} & 11.27 & 0.0 & 0.0210 & 129.0 \\
Tiny-ImageNet & 64 & \texttt{uniform} & 44.23 & 0.0 & 0.0220 & 141.0 \\
Tiny-ImageNet & 64 & \texttt{gap} & 44.52 & 288.0 & 0.0190 & 687.2 \\
\bottomrule
\end{tabular}
\end{table}

\begin{table}[h]\centering\small
\caption{Routing entropy by arm, CIFAR-100 width 64, per-seed means.}
\label{tab:entw64}
\begin{tabular}{l rrr r}
\toprule
arm & mean $\tilde H$ & min & max & $1-\tilde H$ \\
\midrule
\texttt{uniform} & 1.0000 & 1.0000 & 1.0000 & 0.00000 \\
\texttt{gap} & 0.9885 & 0.9867 & 0.9912 & 0.01147 \\
\texttt{shuffled} & 0.9996 & 0.9995 & 0.9997 & 0.00040 \\
\texttt{frozen\_random} & 0.9971 & 0.9967 & 0.9977 & 0.00287 \\
\bottomrule
\end{tabular}
\end{table}

\begin{table}[h]\centering\small
\caption{Every run in the study. GPU-hours are from logged wall-clock
time.}
\label{tab:invent}
\begin{tabular}{l r r l l r r}
\toprule
dataset & width & $K$ & $\tau$ & arms & runs & GPU-h \\
\midrule
CIFAR-100 & 16 & 4 & constant & \texttt{static}, \texttt{uniform}, \texttt{gap} & 18 & 5.5 \\
CIFAR-100 & 32 & 4 & constant & \texttt{static}, \texttt{uniform}, \texttt{gap} & 18 & 6.1 \\
CIFAR-100 & 64 & 4 & constant & \texttt{static}, \texttt{uniform}, \texttt{gap}, \texttt{shuffled}, \texttt{frozen\_random} & 30 & 30.8 \\
CIFAR-100 & 32 & 2 & constant & \texttt{gap} & 6 & 2.8 \\
CIFAR-100 & 32 & 8 & constant & \texttt{gap} & 6 & 6.9 \\
CIFAR-100 & 16 & 4 & annealed & \texttt{gap} & 6 & 4.4 \\
CIFAR-100 & 32 & 4 & annealed & \texttt{gap} & 6 & 3.8 \\
CIFAR-100 & 64 & 4 & annealed & \texttt{gap} & 6 & 8.6 \\
Tiny-ImageNet & 32 & 4 & constant & \texttt{static}, \texttt{uniform}, \texttt{gap} & 18 & 39.5 \\
Tiny-ImageNet & 64 & 4 & constant & \texttt{static}, \texttt{uniform}, \texttt{gap} & 18 & 95.7 \\
\midrule \multicolumn{5}{l}{total} & 132 & 204.1 \\
\bottomrule
\end{tabular}
\end{table}

\end{document}
